\subsection{Обзор аналогов}
\label{sec:analysis:analogues}

Для сопоставления архитектурных и функциональных решений выбраны три профильные системы торговли автомобильными деталями: федеральный интернет-магазин новых компонентов \textit{Exist.ru}, ведущий белорусский портал поиска автозапчастей \textit{Bumper.by} и сетевой интернет-магазин \textit{Autodoc.ru}. Каждый аналог рассмотрен по единой схеме: назначение, принцип функционирования, достоинства и недостатки.

\subsubsection{Специализированный интернет-магазин автозапчастей \textit{Exist.ru}}

\textit{Exist.ru}~-- одна из крупнейших торговых платформ на рынке новых автомобильных запчастей и расходных материалов. Поиск деталей организован по артикулам, \textit{VIN}-коду транспортного средства и структурированным каталогам автопроизводителей с матрицами кросс-кодировок аналогов [1]. Главная страница сайта \textit{Exist.ru} показана на рисунке~\ref{fig:analysis:exist}.

\begin{figure}[H]
	\centering
	\includegraphics[width=0.85\linewidth]{attachments/analogues/exist_screen.png}
	\caption{Главная страница сайта \textit{Exist.ru}}
	\label{fig:analysis:exist}
\end{figure}

К достоинствам системы \textit{Exist.ru} относятся обширная номенклатурная база оригинальных деталей и сертифицированных заменителей, развитые алгоритмы подбора по идентификационному номеру автомобиля и разветвленная сеть пунктов выдачи заказов. Недостатки решения заключаются в закрытом исходном коде, исключающем локальное внедрение для независимого магазина, высокой ресурсоемкости инфраструктуры и отсутствии простых инструментов прямого обмена каталогом в открытых форматах.

\subsubsection{Портал поиска автозапчастей \textit{Bumper.by}}

\textit{Bumper.by}~-- специализированный белорусский интернет-портал, агрегирующий предложения по продаже контрактных, б/у и новых автозапчастей от сотен авторазборок и торговых складов Беларуси [2]. Поиск позиций выполняется по марке, модели, году выпуска автомобиля, типу кузова, узлу и агрегату. Интерфейс поиска автозапчастей сайта \textit{Bumper.by} показан на рисунке~\ref{fig:analysis:bumper}.

\begin{figure}[H]
	\centering
	\includegraphics[width=0.85\linewidth]{attachments/analogues/bumper_screen.png}
	\caption{Интерфейс поиска автозапчастей сайта \textit{Bumper.by}}
	\label{fig:analysis:bumper}
\end{figure}

Достоинствами портала \textit{Bumper.by} выступают крупнейшая база автозапчастей в наличии на складах в Республике Беларусь, публикация реальных фотографий конкретных деталей и прямые контакты продавцов. Ограничения площадки связаны с отсутствием единой онлайн-корзины и сквозного оформления заказов, отсутствием строгого контроля актуальности складских остатков в реальном времени и невозможностью автоматического формирования первичных документов и складской отчетности.

\subsubsection{Интернет-магазин автозапчастей \textit{Autodoc.ru}}

\textit{Autodoc.ru}~-- популярная торговая сеть и интернет-магазин автозапчастей, ориентированный на розничные и оптовые онлайн-продажи [3]. Система обеспечивает регистрацию личного кабинета, подбор деталей, корзину, пошаговое оформление покупки и отслеживание доставки. Каталог деталей сайта \textit{Autodoc.ru} показан на рисунке~\ref{fig:analysis:autodoc}.

\begin{figure}[H]
	\centering
	\includegraphics[width=0.85\linewidth]{attachments/analogues/autodoc_screen.png}
	\caption{Каталог деталей сайта \textit{Autodoc.ru}}
	\label{fig:analysis:autodoc}
\end{figure}

Среди достоинств системы \textit{Autodoc.ru} выделяются продуманный пользовательский интерфейс витрины и корзины, прозрачное информирование покупателя о стадиях транспортировки и развитая система клиентских скидок. В качестве недостатков отмечаются закрытая проприетарная платформа без возможности локальной адаптации, отсутствие открытого реляционного интерфейса к структуре базы данных и нехватка гибкого пакетного обмена каталогами в форматах \textit{CSV}, \textit{JSON} и \textit{XML} для внешних систем.

\subsubsection{Сравнительный анализ систем-аналогов}

Сопоставление рассмотренных систем приведено в таблице~\ref{table:analysis:analogues}.

\begin{table}[ht]
\caption{Сравнительный анализ систем-аналогов}
\label{table:analysis:analogues}
\centering
	\begin{tabular}{|>{\raggedright}m{0.28\textwidth}
	                |>{\raggedright}m{0.21\textwidth}
	                |>{\raggedright}m{0.21\textwidth}
	                |>{\raggedright\arraybackslash}m{0.20\textwidth}|}
	\hline
	\centering Критерий & \centering Exist.ru & \centering Bumper.by & \centering\arraybackslash Autodoc.ru \\
	\hline
	Модель системы & Сетевой интернет-магазин & Агрегатор объявлений & Сетевой интернет-магазин \\
	\hline
	Встроенная корзина и заказ & Реализовано & Отсутствует & Реализовано \\
	\hline
	Транзакционный учет остатков & Реализовано & Отсутствует & Реализовано \\
	\hline
	Пакетный обмен (\textit{CSV}, \textit{JSON}, \textit{XML}) & Ограничено & Отсутствует & Ограничено \\
	\hline
	Аналитические отчеты (\textit{PDF}, \textit{Excel}, \textit{Word}) & Базовые формы & Отсутствует & Базовые накладные \\
	\hline
	Доступность для локального развертывания & Отсутствует & Отсутствует & Отсутствует \\
	\hline
	Мультиязычный интерфейс & Региональный & Только русский & Только русский \\
	\hline
	\end{tabular}
\end{table}

Анализ показывает, что существующие решения либо относятся к закрытым корпоративным платформам, недоступным для локального развертывания независимым предприятием, либо функционируют в формате агрегаторов объявлений без сквозного транзакционного оформления заказов и автоматизированной складской отчетности.

Формируется свободная ниша для компактной и производительной информационной системы, сочетающей функции интернет-магазина, транзакционный контроль складских остатков, обмен каталогами в открытых форматах (\textit{CSV}, \textit{JSON}, \textit{XML}) и автоматическую генерацию аналитических отчетов без избыточных накладных расходов.
